\begin{algorithm}[H]
\caption{Exact maximum depth with a distinguished miss hypothesis and bounded caps}
\label{alg:bounded-caps}
\KwIn{Distinguished hypothesis $p$; other hypotheses $P$; supplied caps of size at most fixed $b$.}
\KwOut{An exact minimum-height supplied-query tree.}
\If{$P=\varnothing$}{\KwRet the zero-query tree\;}
Compute the minimum supplied-halfplane cover size $K$ of $P$\;
\For{$D=K,\ldots,K+b-1$}{
  $d_0\gets\max\{0,D-b+1\}$\;
  \ForEach{ordered supplied-query suffix of length $q\le D-d_0$}{
    \ForEach{$R$ contained in the union of its caps}{
      \If{the suffix does not cover $R$}{\Continue\;}
      Find a minimum cover $Q$ of $P\setminus R$ using only caps disjoint from $R$\;
      \If{$|Q|>d_0$ or the cover is infeasible}{\Continue\;}
      Compute the suffix's first-hit groups $A_j$ on $R\cup\{p\}$\;
      \If{some $j+c_W(\ExactTree(A_j,\mathrm W))>D-d_0$}{\Continue\;}
      Execute $Q$ followed by the suffix on the all-miss spine\;
      Attach exact local supplied-query trees to all positive first-hit groups\;
      \KwRet this tree\;
    }
  }
}
\end{algorithm}

The theorem's height interval means the loop needs only the listed
values of $D$. For each value, the algorithm searches suffixes and
residual sets; since a suffix has at most $b-1$ queries and each cap
contains at most $b$ points, its union contains at most $b(b-1)$
points. An empty suffix therefore permits only $R=\varnothing$. The
prefix step is a minimum supplied-halfplane cover, computed by the
polynomial algorithm cited in Theorem~\ref{thm:bounded-cap-height};
when the instance has the special hull-interval representation, its
dedicated interval oracle performs the same step.
